\enabstractparagraph{With the growth of laboratory equipment in universities, lending registration, return confirmation, and overdue reminders have become frequent management tasks. Paper forms and isolated spreadsheets make it difficult to keep inventory states consistent. This thesis designs and implements a campus laboratory equipment lending management system for students and laboratory administrators.}

\enabstractparagraph{The system adopts a front-end and back-end separated architecture. Equipment records, lending requests, approval records, return confirmation, and audit logs are organized as independent modules. Students can search available equipment, submit requests, and track progress. Administrators can maintain equipment status, approve requests, register returns, and inspect overdue records. A state machine constrains the request workflow, while transactions keep lending quantities and equipment inventory synchronized.}

\enabstractparagraph{The work covers requirements analysis, system design, database design, implementation, and testing. Functional tests include login, equipment search, request approval, return registration, and invalid input handling. The results show that the example system completes the lending workflow as designed and records audit information for critical operations.}

\enkeywords{laboratory management; equipment lending; state machine; transaction processing; audit log}
